\section{Identity, discovery and mission association}\label{sec:discovery}
\subsection{Identity model}
\begin{longtable}{L{34mm}L{127mm}}
\toprule Field & Definition\\\midrule
\endfirsthead
\toprule Field & Definition\\\midrule
\endhead
\code{node_id} & Persistent UUID for an application installation; independent of hostname, IP and chosen role. Cloned installations must generate distinct IDs.\\
\code{boot_id} & Fresh UUID on every service process start. Sequences and monotonic times are meaningful only within this boot.\\
\code{station} & Exactly \code{launch}, \code{gs1}, \code{gs2}, \code{gs3} or \code{gs4}.\\
\code{role} & Exactly \code{telem} or \code{video}. Together with station, identifies a logical slot.\\
\code{mission_id} & UUID shared deliberately across the ten computers; \code{null} until assigned. A human-readable mission name is not a unique identity.\\
\code{vehicle_mode} & \code{balius} or \code{iris}; must agree with the mission.\\
\code{protocol} & Object containing integer \code{major} and \code{minor}; this draft specifies 1.0. App version is a separate string.\\
\bottomrule
\end{longtable}

\req{ID-01}{Two live nodes claiming the same slot SHALL create a visible
conflict. Neither SHALL replace the existing selected node automatically.
DNS name conflict resolution does not resolve logical slot conflicts. Concurrent
endpoints presenting the same node ID with different live boots SHALL also be
flagged as an identity conflict rather than merged.}
\req{ID-02}{The wizard's remember checkbox MAY persist user choices; permanent
node identity is separate. Changing station, role, mission or source mode while
serving SHALL end current subscriptions/media sessions and start a new boot
identity. Network data from the previous identity SHALL not leak into the new
session. No sequence reset is permitted inside one boot.}

Mission assignment is explicit: Launch Telemetry creates or opens a mission
UUID; operators join it in the other applications or supply it on the CLI.
A receiving UI can display a mismatched peer in diagnostics, but SHALL exclude
its data from that mission's telemetry/video workspace. If Launch Telemetry is
absent, a previously shared mission UUID still works. There is no implicit
assignment from an SSID, a Bonjour name, or identical text such as ``Iris Launch.''

\subsection{DNS-SD records and names}
The proposed private service type is \code{_rt-gs._tcp.local.}; this document
does not claim an IANA service-name registration. An example instance is
\code{rt-gs1-telem._rt-gs._tcp.local.}. PTR browsing finds instances;
SRV supplies the HTTP host/port; TXT supplies compact hints. These are standard
DNS-SD mechanisms \cite{dnssd}. The app SHALL publish after its API is listening
and withdraw the service on clean shutdown.

\begin{lstlisting}
Service type: _rt-gs._tcp.local.
Instance:     rt-gs1-telem
SRV target:   rt-gs1-telem.local.
SRV port:     8766
TXT: p=1; node=<UUID>; station=gs1; role=telem; api=/api/v1
\end{lstlisting}
Each TXT key/value is a separate DNS TXT string; the semicolon layout above is
display notation only. Total TXT payload SHALL be at most 512 bytes in this
profile. No credentials, telemetry samples or rapidly changing status belong
there. The HTTP identity response is authoritative; TXT is an untrusted discovery
hint. A disagreement SHALL block normal attachment until rediscovery resolves it.

Preferred hostnames follow the user's \code{rt-base-*}/\code{rt-gsN-*}
convention. If an application cannot safely publish that hostname, SRV MAY target
the machine's existing resolvable hostname or app-owned mDNS host name. Clients
SHALL use SRV/A/AAAA resolution and SHALL not construct a URL by appending
\code{.local} to the service instance. Global OS hostname changes are unnecessary.

mDNS operates on the local link using UDP 5353; cached records are not heartbeats
\cite{mdns}. Access points must permit peer traffic and multicast. Routed VLANs
require an explicitly configured discovery gateway or manual endpoints. The
protocol SHALL support direct unicast IPv4 and IPv6 endpoints; link-local IPv6
addresses require the local interface scope. URLs use bracketed IPv6 literals.

\subsection{Ports and compatibility}
\begin{longtable}{L{35mm}L{34mm}L{92mm}}
\toprule Function & Proposed default & Rule\\\midrule
\endfirsthead
\toprule Function & Proposed default & Rule\\\midrule
\endhead
mDNS & UDP 5353 & Managed by the discovery implementation.\\
HTTP + WebSocket & TCP 8766 & Configurable; advertise the actual bound port. Do not silently bind a different service at the requested port.\\
SRT video sessions & UDP 50000--50063 & Configurable listener allocation range on video nodes; return actual session endpoints. Client source ports are OS-assigned.\\
Old snapshot link & TCP 8765 & Existing v0a protocol is distinct. No wire compatibility or automatic bridge is implied.\\
\bottomrule
\end{longtable}

\req{DISC-01}{Automatic discovery SHALL fall back to an operator-entered
HTTP endpoint. Manual endpoints undergo the same identity/mission/version checks.
An IP change SHALL trigger re-resolution and a new handshake, not a new logical
station. Duplicate appearances of the same node SHALL merge only when boot and
identity agree.}
\req{DISC-02}{Major-version mismatch SHALL be incompatible. For major 1,
unknown optional fields SHALL be ignored and missing required fields rejected.
New required behavior needs capability negotiation or a new major version.
An unsupported event kind SHALL be logged and ignored without losing its cursor.}

\section{HTTP API and common encoding}\label{sec:http}
The normative monitoring baseline uses HTTP/1.1 and WebSocket on a team-controlled
LAN. It provides no cryptographic peer authentication or confidentiality. A node
UUID and discovery name are identifiers, not proof of trust. Network publication
SHALL be visible to the operator. Expose no arbitrary files or shell execution.
HTTPS/WSS with provisioned identities is a future compatible transport profile;
it is required to be designed before enabling remote hardware commands on an
untrusted network. This draft's network interface cannot actuate flight hardware.

All JSON uses UTF-8 and \code{application/json}; no NaN or infinity is legal.
Unavailable measurements use \code{null}, never zero or an empty string.
Counters that can grow beyond JavaScript's exact integer range use decimal
strings with no leading zeros except \code{"0"}. UUIDs use canonical lower-case
hyphenated form. UTC timestamps use RFC 3339 with \code{Z}; \code{null} is allowed
when wall-clock time is unavailable. Monotonic milliseconds are decimal strings.
Request bodies are limited to 64 KiB; JSON responses and WS messages to 256 KiB.
Raw packets within a message are limited to 64 KiB decoded bytes. Larger future
payloads require another negotiated profile, not silent truncation.

\begin{longtable}{L{19mm}L{61mm}L{80mm}}
\toprule Method & Path under \code{/api/v1} & Semantics\\\midrule
\endfirsthead
\toprule Method & Path under \code{/api/v1} & Semantics\\\midrule
\endhead
GET & \code{/info} & Identity, capabilities, clocks and protocol version.\\
GET & \code{/status} & Atomic current hardware/network state with its event watermark.\\
GET & \code{/telemetry/latest} & Latest accepted observation per known vehicle; unknown attribution is a separate entry. Empty array on video-only nodes.\\
GET & \code{/streams} & Channel/rendition manifest; empty array on telemetry-only nodes.\\
GET (WS) & \code{/events} & WS event subscription; subprotocol \code{rt-gs.v1}.\\
POST & \code{/video/sessions} & Create one authorized-by-profile receive session, detailed in Section~\ref{sec:video}.\\
PATCH & \code{/video/sessions/\{id\}} & Renew that session's lease.\\
DELETE & \code{/video/sessions/\{id\}} & Release its resources; repeat is harmless.\\
POST & \code{/commands} & Reserved; return 501 \code{unsupported} in this revision. No dispatch.\\
\bottomrule
\end{longtable}

GET responses SHALL include \code{Cache-Control: no-store}. No telemetry endpoint
is backed by a constantly rewritten \code{telem.txt}. Clients SHOULD use a
five-second HTTP timeout and cancel obsolete requests after identity changes.
The server SHALL reject non-local browser origins by default rather than provide
wildcard cross-origin access; this is not a substitute for authentication.

Every successful GET body includes \code{protocol}, \code{node_id},
\code{boot_id}, \code{mission_id} and \code{generated_mono_ms}, followed by its
endpoint payload. Snapshot endpoints also include \code{as_of_event_seq}, the
atomic journal watermark. A health response is independent of packet freshness.
\code{/info} additionally contains station, role, vehicle mode, app version,
source mode, supported event kinds, resource limits and capabilities. Capability
Booleans describe actual implementation; a role label is not a capability.

Error bodies contain \code{error.code}, \code{error.message},
\code{error.retryable} and optional \code{request_id}. HTTP status mapping:
400 malformed JSON; 404 unknown resource; 405 method; 409 identity/mission or
idempotency conflict; 410 expired session; 413 size; 415 content type; 422 invalid
field or channel; 429 capacity (with \code{Retry-After}); 501 unsupported operation;
503 device/pipeline unavailable. Errors SHALL not expose secrets or stack traces.
GET never changes hardware state.
